\chapter{酶：让正确反应在正确时间发生}

\begin{kaichang}
超市调料区有一种白色粉末，叫“嫩肉粉”。用法简单得近乎可疑：撒一点在又老又韧的生牛肉上，拌匀，常温下放半小时——不用加热，不用加醋，什么都不用做。半小时后下锅，原本嚼不动的肉变得又滑又嫩。

请猜一猜：这包粉末对肉做了什么？是把它“泡软”了？还是偷偷加热了它？还是有什么东西在“咬”这块肉？再看看包装袋上的小字：主要成分写着“木瓜蛋白酶”，注意事项里还叮嘱——不要用开水冲调。为什么怕开水？

把答案先押在纸上。这一章结束时，我们回来开奖。而这包不起眼的粉末，将带我们跨过整本书最重要的一条边界：从“物质的化学”跨进“生命的化学”。
\end{kaichang}

\section{厨房里的三桩怪事}

嫩肉粉不是唯一的怪事。再给你两个，都是我敢保证你见过的。

第一桩：把一块没滋没味的白馒头放进嘴里，慢慢嚼，别急着咽。嚼上一分钟，你会尝到一股明显的甜味。馒头配料表里可没有糖。

第二桩：处理伤口用的双氧水，静置在瓶子里几年都不怎么变样；可只要滴一小滴在切开的生猪肝（或者生土豆）上，立刻像烧开了一样冒出大量气泡——那是氧气。没有加热，没有通电，只是碰到了一点肝。

第三桩：沾了奶渍、血渍的衣服，普通肥皂搓不干净，换“加酶洗衣粉”泡一泡就干净了。但说明书又叮嘱：用 $40\,^{\circ}\mathrm{C}$ 左右的温水效果最好，别用滚烫的水。

这三桩怪事看起来毫不相干：肉、馒头、双氧水、衣服。但把它们摆在一起，共性就浮出来了：

\begin{itemize}[itemsep=2pt]
\item 变化都在\textbf{温和条件}下发生——常温、常压、不用强酸强碱；
\item 只需要\textbf{极少量的“神秘成分”}——一勺粉末、一口唾液、一滴肝汁；
\item 每种神秘成分\textbf{只管一件事}——嫩肉粉不会去分解奶渍，唾液不会让双氧水冒泡；
\item 它们都\textbf{怕烫}。
\end{itemize}

一百多年前的化学家也被同样的怪事困扰着。他们已经能在烧瓶里催化各种反应（第 30 章讲过的铂、铁、硫酸），可实验室的催化剂要么需要高温高压，要么“什么反应都催、但什么都不精”。而生物体内这些神秘成分，专一、高效、温和得不像话。它们后来被统称为\textbf{酶}——一种由生物制造的催化剂。

\section{推进镜头：一只折叠出来的“口袋”}

第 48 章讲过，蛋白质是一条氨基酸长链折叠成的三维结构，折叠方式决定了它能干什么。酶，就是一类专门折叠出“工作间”的蛋白质。

这个工作间叫\textbf{活性位点}：蛋白质表面一处形状特殊的小口袋或小沟，通常只占整个分子表面的很小一块。口袋里排列着几个精心挑选的氨基酸侧链——有的带正电，有的带负电，有的疏水，有的能形成氢键（第 48 章讲过的那些侧链性格，在这里全用上了）。能被这只口袋接纳、并在里面发生反应的分子，叫做这个酶的\textbf{底物}。

底物分子在溶液里乱撞，一旦撞进活性位点、形状和电荷恰好对得上，就会被“抓”住：侧链们用氢键、静电吸引、疏水作用把它固定在特定的朝向和位置上。这个临时组合叫\textbf{酶—底物复合物}。随后，口袋里的侧链开始“动手”：有的拉扯底物里某根化学键，有的递过去一个质子，有的暂时接过一对电子。反应完成后，产物与口袋不再匹配，脱落离开——酶恢复原样，等待下一个底物。

最早描述这种专一性的是 1894 年的德国化学家费歇尔（就是第 40 章酯化反应提到的那位）。他提出：酶和底物的关系像\textbf{锁和钥匙}——口袋是固定的锁孔，只有形状吻合的钥匙插得进去。这个比喻解释了专一性：唾液淀粉酶的锁孔只认淀粉，双氧水的钥匙插不进去。

但锁钥说有个破绽：如果锁孔的形状分毫不能变，底物进去时的形状和反应中途的形状并不一样，锁怎么能同时配合两者？1958 年，美国科学家科什兰提出了更贴近真实的图景——\textbf{诱导契合}：活性位点不是硬邦邦的锁孔，而更像一只\textbf{手套}。手伸进去之前，手套是瘪的；手伸进去时，手套被撑开、裹紧，正好贴合手的形状。同样，底物进入口袋时，活性位点的形状会微微调整，把底物裹得更紧——这一“裹”，恰好把底物挤进了最有利于反应的姿态：该被拉长的键被拉长，该靠近的两个基团被凑到一起，该屏蔽的水分子被挡在外面。

也就是说，酶不只是“认出”底物，它还\textbf{用力改变了底物的处境}。这句“用力”，就是下一节的主角。

\section{酶真正做的事：把山坡削矮}

回忆第 29 章的图景：一个反应哪怕放能、哪怕“愿意”发生，分子也得先翻过一座\textbf{活化能}的山坡——反应物要先被掰弯、拉长、凑近到一个别扭的过渡状态，才能滑向产物。山坡越高，能翻过去的分子越少，反应越慢。第 30 章还讲过催化剂的通用玩法：它不推分子，而是\textbf{修一条翻山的隧道}——提供另一条活化能更低的反应路径。

酶干的就是这件事，只是修隧道的手艺登峰造极。诱导契合那一“裹”，把过渡状态稳稳托住：底物被摆成接近过渡态的姿势，口袋里的电荷分布又恰好能安抚过渡态上的电荷——相当于把山尖削掉了一大截。

下面这张图是全章最该记住的一幅（这是人为的表示法，横轴“反应进程”不是真实距离）：

\begin{center}
\begin{tikzpicture}[scale=0.9]
\draw[->] (0,0) -- (10.5,0) node[right]{反应进程};
\draw[->] (0,0) -- (0,4.8) node[above]{能量};
\draw[thick] (0.5,1.2) .. controls (3,5.2) and (4,5.2) .. (6,2.6)
   .. controls (7,1.8) .. (9.5,0.7);
\draw[thick,dashed] (0.5,1.2) .. controls (2.5,3.2) and (3.5,3.2) .. (5.5,2.2)
   .. controls (7.2,1.5) .. (9.5,0.7);
\node at (1.2,0.7) {反应物};
\node at (9,0.2) {产物};
\node at (4.6,4.6) {无催化};
\node at (2.9,2.6) {有酶};
\draw[<->] (8.2,0.7) -- (8.2,1.2) node[midway,right]{$\Delta G$};
\end{tikzpicture}
\end{center}

实线是没有酶时的山坡，虚线是有酶时修出的矮坡。请注意两件事。第一，\textbf{起点和终点的高度没有变}：反应物和产物的能量差 $\Delta G$ 由分子本身决定（第 28 章），酶碰都没碰它。第二，由此推出一条经常被考倒人的结论：\textbf{酶不改变平衡}。第 31 章讲过，平衡位置由 $\Delta G$ 说了算；既然 $\Delta G$ 不变，加不加酶，最终“反应物和产物各占多少”完全一样。酶只是让你\textbf{更快到达}那个终点——而且因为它同时降低正、逆两个方向的活化能，两个方向一起变快，谁也不偏袒。

所以请记住酶的三条“不为”：\textbf{不提供能量}（它不推分子，只削矮山坡）；\textbf{不改变方向}（$\Delta G$ 为正、本来不愿发生的反应，加再多酶也不会发生）；\textbf{不改变终点}（平衡位置纹丝不动）。这三条后面还要反复用。

那酶把山坡削矮了多少？用数字感受一下。

\begin{qiaoliang}
以双氧水分解为例：\chem{H_2O_2} 分解成水和氧气，没有催化剂时慢到一瓶双氧水能放好几年。猪肝里的\textbf{过氧化氢酶}把它削成了平路：\textbf{一个}过氧化氢酶分子，每秒钟大约能分解 $10^{5}$ 到 $10^{6}$ 个双氧水分子——是自然界最快的酶之一。算一笔账：一滴肝研磨液里大约有 $10^{15}$ 个酶分子，全部开工每秒能处理约 $10^{20}$ 个双氧水分子。一滴 $3\%$ 的双氧水里有多少分子？大约 $10^{19}$ 个。也就是说，一滴肝汁碰到一滴双氧水，理论上\textbf{不到一秒}就能把对方的分子清空——这就是你看见气泡“瞬间沸腾”的原因。而没有酶时，同样这批分子要靠运气慢慢翻山，得等上好几年。同一座山，有隧道和没隧道，时间差了约 $10^{8}$ 倍。许多酶催化反应的加速倍数在 $10^{5}$ 到 $10^{17}$ 之间——没有酶，你身体里绝大多数反应慢得跟没有一样，“活着”这件事在化学上根本不可能。
\end{qiaoliang}

\section{符号登场：\texorpdfstring{$\chem{E}$ 加 $\chem{S}$}{E 加 S}}

规律看清了，现在才轮到符号。酶催化的通用写法是：

\[\chem{E} + \chem{S} \rightleftharpoons \chem{ES} \rightarrow \chem{E} + \chem{P}\]

$\chem{E}$ 是酶，$\chem{S}$ 是底物，$\chem{ES}$ 是酶—底物复合物，$\chem{P}$ 是产物。注意酶在左边进去、右边原样出来——和第 30 章的所有催化剂一样，它\textbf{不被消耗}。这解释了一勺嫩肉粉为什么能嫩一大块肉：每个酶分子可以反复“接单”成千上万次。

唾液让馒头变甜的反应，写出来是淀粉的水解（第 46 章细讲过糖和糖苷键）：

\[\chem{(C_6H_{10}O_5)_{n}} + n\chem{H_2O} \rightarrow n\chem{C_6H_{12}O_6}\]

淀粉长链在唾液淀粉酶的口袋里被逐段剪开，先变成短链糊精和麦芽糖——你尝到的那点甜，就是被剪下来的小糖。没有这个酶，同样这口馒头要在肚子里泡上不知多久才能派上用场。

还有一个符号层面的规律值得展开：反应速率怎样随底物浓度变化。把酶的用量固定，逐渐增加底物，速率会先几乎成比例地上升，然后越来越平，最后停在一条水平线上——像一根被拉平的增长曲线。为什么？低浓度时，大量酶的口袋空着，多来些底物就有更多口袋在干活；可底物多到某个程度，所有口袋\textbf{全天候满员}，酶分子连轴转也接不过来了，再加底物只能排队——速率达到上限。这个上限记作 $V_{\max}$（最大速率），达到一半 $V_{\max}$ 时的底物浓度记作 $K_{\mathrm{m}}$（米氏常数），它可以粗略理解为“口袋对底物的亲和力指标”：$K_{\mathrm{m}}$ 越小，越少的底物就能让酶半饱，抓得越紧。

\begin{shiyan}
亲手验证你的唾液里有酶（安全可在家做）。材料：一小块馒头或一勺米饭、两个干净小碟、药店买的碘伏（碘液）、温水、一根干净筷子。步骤：①把馒头撕成两份差不多大的小团；②第一份放进嘴里充分咀嚼一分钟，让唾液浸透，吐到碟子里（自己的唾液自己用，不与他人混用）；③第二份加等量的温水，用筷子同样搅成糊状；④两份各滴一滴碘液，观察颜色。观察点：第 46 章讲过，淀粉遇碘变深蓝。温水那份应当明显变蓝，而嚼过的那份蓝色明显变浅甚至不显色——因为淀粉已经被唾液淀粉酶剪碎了一部分。安全提示：碘液只许外用、切勿入口，实验后洗手；所有接触唾液的用具自己专用并洗净。想再进一步，可以把两份样品同时放进约 $37\,^{\circ}\mathrm{C}$ 的温水浴里保温五分钟再滴碘，对比更清楚——顺便想想，为什么接近体温效果最好？
\end{shiyan}

\section{科学家怎样知道：酶不是“生命力”}

\begin{zhengju}
19 世纪，酿酒和发面早已是千年老手艺，但“发酵到底是什么”吵了几十年。以巴斯德为代表的一派坚持：发酵离不开\textbf{活的}酵母细胞，糖变酒精靠的是某种只存在于活细胞里的“生命力”——化学反应本身做不到。这话听起来有道理：毕竟你把酵母煮死，发酵就停了。

1897 年，德国化学家爱德华·毕希纳本想用酵母榨汁做点药物研究，顺手往榨出的酵母汁里加了糖防腐——结果汁里冒泡了：糖在无细胞的发酵汁里变成了酒精。\textbf{没有活细胞，发酵照样发生}。所谓“生命力”，其实是酵母细胞里一群可以离开细胞、独立干活的物质——酶。毕希纳的实验设计其实朴素得可爱：活细胞组与无细胞汁组，其他条件相同，只动一个变量。替代解释（“可能还有没榨碎的活细胞”）也被显微镜检查和反复稀释排除了。这一罐冒泡的酵母汁，把“生命化学”从神秘主义手里夺回来，交给了化学。

下一个问题接踵而至：酶这种“物质”到底是什么？1926 年，美国化学家萨姆纳从刀豆里提纯出一种酶——脲酶，把它\textbf{结晶}了出来。在当时，结晶是纯净化合物的标志，而晶体的成分是蛋白质。萨姆纳宣布“酶就是蛋白质”，迎来的却是冷嘲热讽——大权威们认为蛋白质太“脏”，不过是吸附了真正活性成分的杂质。争论持续了近十年，直到更多酶被一个个结晶出来、个个都是蛋白质，共识才翻转。后来又发现极少数具有催化能力的 RNA（核酶），“酶全是蛋白质”这句话才补上了一个小脚注——科学结论就是这样，一次次被证据逼着打补丁。

至于“酶和底物到底怎样结合”，就是前面讲过的故事：1894 年费歇尔的锁钥模型统治了半个多世纪，1958 年科什兰用诱导契合修正了它——不是推翻，而是把“锁孔会动”这一层加了进去。这个修正有实验支撑：人们后来用 X 射线晶体学直接拍到同一个酶“没结合底物”和“结合了底物”的两张照片，形状确实不一样。从“生命力”到“蛋白质”到“会变形的口袋”，每一步都不是天才的灵光，而是\textbf{旧解释被实验顶到墙角之后的让步}。
\end{zhengju}

\section{酶的软肋：它会坏、会被堵、会被算错}

讲到这里，酶近乎完美。但正因为它是一台依赖精密折叠结构的蛋白质机器（第 48 章），它有几条实实在在的软肋。

\textbf{第一，它会被“烫散”。} 温度升高，分子运动加剧，反应通常变快——所以酶促反应的速率起初随温度上升。但超过某个温度，维持折叠的氢键和疏水作用开始撑不住，蛋白质散开（第 48 章讲过的变性），活性位点的口袋塌了，酶就永久下岗。这就是嫩肉粉和洗衣粉都怕开水的原因。对大多数来自人体的酶，最舒适的温度就在体温附近；温泉细菌里的酶却能扛到 $70\,^{\circ}\mathrm{C}$ 以上——酶的“适宜温度”取决于它出身的生物，没有统一标准。

\textbf{第二，它挑酸碱度。} 口袋里的侧链带不带电，取决于环境 pH（第 32 章）；pH 一变，口袋的电荷和形状就跟着变。胃里的胃蛋白酶在 pH 约 $2$ 的强酸里最活跃，到了小肠的弱碱性环境就蔫了；接班的是胰蛋白酶，恰好喜欢 pH 约 $8$。消化道一路走下来，pH 像接力棒一样换，酶也一批批换岗。

\textbf{第三，它会被人“堵”——这恰恰是许多药物的工作原理。} 想一想：既然酶靠口袋认底物，那么塞一个形状相似、却不能正常反应的分子进去，口袋就被占住了。这叫\textbf{竞争性抑制}：抑制剂和底物抢同一个座位，底物浓度足够大时还能把座位抢回来。最早的抗菌药磺胺就是这个思路：细菌需要一种酶来制造叶酸，磺胺分子长得和这种酶的底物几乎一样，抢座位占着不干活，细菌断供叶酸、停止繁殖；人体细胞不自己造叶酸（我们从食物里获取），所以基本不受误伤——好药物的本质，就是找到\textbf{敌人有、你没有}的那只口袋。另一类堵法更彻底：阿司匹林（第 44 章）会把一个叫环氧酶的酶的活性位点\textbf{永久焊死}（共价结合），这个酶分子从此报废，止痛效果要等身体重新合成新的酶才消退——这叫不可逆抑制。还有一类不抢座位的：抑制剂结合在口袋\textbf{之外}的位置，让整个蛋白质变形、口袋跟着走样，底物照常来却待不住——这叫\textbf{非竞争性抑制}，此时加再多底物也救不回速率。许多毒物（比如某些重金属离子）就是这样让关键酶集体变形的。

\textbf{第四，关于它的漂亮公式，只在限定条件下成立。} 前面“速率先升后平”的规律，1913 年被米凯利斯和门滕写成了一个简洁的方程（见章末挑战层），它成功描述了许许多多酶。但它成立的前提是几条假设：测的是反应\textbf{刚开始}的速率（产物还没积累、逆反应可忽略）；底物远多于酶；酶—底物复合物浓度近似不变；而且酶只有\textbf{一种}底物、不受调节。细胞里的真实情况常常不满足——有的酶受产物“反馈”抑制，有的酶被别的分子远程调节，有的底物本身会改变酶的脾气（血红蛋白那样的变构效应）。公式没错，但把它当成任何情况下的真理，就是模型的越界——这正是“边界”二字的含义。

\begin{wuqu}
“喝酵素口服液，能补充身体里的酶，助消化、排毒、抗衰老。”——这类产品卖得火热，但它在两个层次上站不住。第一，酶是蛋白质，你喝下去的“酵素”进入胃里，首先遇到的是胃酸和胃蛋白酶——一台专门拆蛋白质的机器。外来的酶会被当成普通蛋白质剪成小段氨基酸再吸收，\textbf{几乎不可能}以完整形态溜进血液去“上岗”。从这个意义上说，喝酶和喝蛋花汤的差别不大。第二，就算有酶侥幸完整吸收，它也认不出该去哪只口袋——酶的专一性意味着它只管自己那一种反应，不存在“万能排毒酶”。需要提醒的是：确实有少数酶制剂可以口服起效，比如乳糖不耐受者（第 46 章）随餐服用的乳糖酶——但那是因为它的工作场所就是肠道内部，拆完乳糖就地退休，根本不需要被吸收。“有效”的边界永远是\textbf{机制},不是广告词。买不买是你的自由，但至少要清楚钱花在了哪个层次。
\end{wuqu}

\section{回到那勺嫩肉粉}

现在可以开奖了。嫩肉粉里的木瓜蛋白酶，是从木瓜汁液里提取的一种蛋白质。撒在肉上，它的活性位点开始识别肉里长长的胶原蛋白和肌肉纤维蛋白，把肽键一根根剪断（第 48 章）——肉的“韧性”正来自这些纠缠的长蛋白链，长链被剪成小段，肉的口感自然就嫩了。

回过头核对开头的每个疑点：

\begin{itemize}[itemsep=2pt]
\item \textbf{为什么常温就行？} 酶把活化能的山坡削矮了，常温下分子自身的热运动就足够翻过矮坡，不用加热。
\item \textbf{为什么只用一点？} 酶不被消耗，一个分子能反复接单成千上万次。
\item \textbf{为什么怕开水？} 高温使蛋白质变性，口袋塌了，隧道塌了——嫩肉粉变成一勺普通蛋白粉。
\item \textbf{为什么它不去“嫩”别的？} 木瓜蛋白酶的口袋认蛋白质，所以你不用担心它把盘子腐蚀了。
\end{itemize}

顺便回应另外两桩怪事：馒头变甜，是唾液淀粉酶在体温附近的口腔里把淀粉剪成了小糖；双氧水冒泡，是猪肝细胞里的过氧化氢酶在飞速解毒——双氧水对细胞有毒，这种酶的本职就是把它迅速拆成水和氧气，你只是借用了它的加班费。

\section{快，还不够：谁来决定“该不该反应”}

走出厨房，酶早已是一支产业大军。加酶洗衣粉里配了蛋白酶、脂肪酶、淀粉酶的“三人组”，分别对付血渍奶渍、油渍和食物残渣——温水让它们处于最佳工况，所以温水洗才干净。食品工业用葡萄糖异构酶把玉米糖浆里的葡萄糖部分转变成更甜的果糖，年产数百万吨高果糖浆；做奶酪靠凝乳酶剪断牛奶里的酪蛋白；制生物燃料要靠纤维素酶拆开植物细胞壁里顽固的纤维素（第 46 章）。医药就更不用说了：据估计，市面上相当大比例的药物，靶点都是某个酶——用抑制剂堵住一只关键的口袋，就是现代药物设计的主流剧本之一。

但我要在这一章结尾泼一盆冷水，也是埋一个伏笔。本章反复强调：酶只决定\textbf{快慢}，不决定\textbf{方向}。可看看你的细胞此刻在干什么——它在把氨基酸拼成蛋白质（吸能）、把小分子拼成大分子（吸能）、把物质从低浓度泵到高浓度（吸能）。这些反应 $\Delta G$ 为正，按第 28 章的判决，它们“不愿意”发生；而酶再神通广大，也动摇不了 $\Delta G$ 一根毫毛。

那生命是怎么让这些“不愿意”的反应天天发生、分秒不停的？答案是：细胞发明了一套“能量转账”系统，把“愿意发生”的放能反应和“不愿意发生”的吸能反应\textbf{捆绑}在一起，用前者拖着后者走。这套转账系统的通用货币，就是你多半听过名字的 ATP。能量到底怎样转账？为什么“ATP 里藏着高能键”这个说法其实有问题？——第 50 章，我们去查细胞的账本。

\begin{tiaozhan}
米凯利斯—门滕方程长这样（跳过不影响主线）：

\[v = \frac{V_{\max}\,[\chem{S}]}{K_{\mathrm{m}} + [\chem{S}]}\]

其中 $v$ 是反应初速率，$[\chem{S}]$ 是底物浓度。自己动手验证它的三个极限：①当 $[\chem{S}]$ 远大于 $K_{\mathrm{m}}$，分母近似等于 $[\chem{S}]$，于是 $v\approx V_{\max}$——口袋全满，速率封顶；②当 $[\chem{S}]=K_{\mathrm{m}}$，$v = V_{\max}/2$——这正是 $K_{\mathrm{m}}$ 的定义；③当 $[\chem{S}]$ 远小于 $K_{\mathrm{m}}$，$v\approx (V_{\max}/K_{\mathrm{m}})\,[\chem{S}]$——速率与底物浓度成正比，口袋大量空置。这个方程是从“$\chem{E}+\chem{S}\rightleftharpoons\chem{ES}\rightarrow\chem{E}+\chem{P}$”的机理出发，假设 $\chem{ES}$ 的浓度在反应初期近似不变（稳态假设）推导出来的。所以一旦底物会被产物抑制、酶受远程调节、或逆反应不可忽略，方程就得改写。一个公式和它的假设是一起出厂的——记住这一点，比记住公式本身更重要。
\end{tiaozhan}

\section*{练一练}

\begin{sikaoti}
\item 画两幅能量图（参照本章的图）：同一个放能反应，一幅无酶、一幅有酶。在两幅图上分别标出活化能箭头和 $\Delta G$ 箭头，然后回答：哪根箭头被酶改变了？哪根没有？由此说明“酶不改变平衡”。
\item 过氧化氢酶一个分子每秒分解约 $10^{6}$ 个双氧水分子。一块猪肝里的这种酶如果全部开工，处理一瓶约含 $10^{22}$ 个双氧水分子、浓度 $3\%$ 的双氧水，需要大约 $10^{16}$ 个酶分子。一块重 \ud{500}{g} 的肝大约有 $10^{11}$ 个细胞，假设每个细胞含 $10^{7}$ 个过氧化氢酶分子，估算一下总酶分子数，说说它够不够、富余多少倍。
\item 做馒头时，面团要在温暖处“醒发”，但蒸馒头时酶却不再起作用。用“温度—速率”曲线（先升后降）解释这两个阶段，并在你画的曲线上标出“醒发温度”和“蒸制温度”大致落在哪一段。
\item 一种竞争性抑制剂和底物结构相似。画一幅“粒子图”：用圆圈表示酶（带缺口表示口袋）、三角表示底物、方块表示抑制剂，画出“只有底物”“加入抑制剂”“底物大大增加”三种情况下口袋的占用情况，并解释为什么竞争性抑制可以被“加底物”逆转，而非竞争性抑制不能。
\item 胃蛋白酶（最适 pH 约 $2$）和胰蛋白酶（最适 pH 约 $8$）如果互换工作环境，各自会怎样？从“侧链带电状态改变口袋形状”的角度给出粒子层面的解释，并说说这对“口服酶类保健品”的设计意味着什么。
\item 有人声称：“给反应体系加一种超级酶，就能把反应物 $100\%$ 全部转化成产物。”用本章的三条“酶之不为”逐条反驳他，并说明：如果真想提高平衡时产物的比例，正确的思路应该动哪个量（提示：回到第 28、31 章）。
\end{sikaoti}
